%قبل از شروع مسئله ترجیحاً چیزی ننویسید. اما اگر خواستید بنویسید، در قالب لیستی (فقرات و شمارش) بنویسید.
%برای اضافه کردن نمره، بعدی از شروع‌مسئله، [x] را اضافه کنید و به جای x نمره مورد نظر را قرار دهید. برای زیرمسئله‌ها هم به همین صورت است.
%\شروع‌مسئله[10]
\شروع‌مسئله [12]
\\
شنگدباو دو تا زیر مجموعه از اعداد $1$ تا $M$ دارد. آن‌ها را $A$ و $B$ بنامید.

مجموعه تمامی اعدادی از $1$ تا $M$ که به صورت جمع دو عضو از $A$ و $B$ قابل نمایش‌هستند را به صورت زیر تعریف می‌کنیم:
$$A * B = \{a + b | a \in A , b \in B, 1 \le a + b \le M\}$$
و به صورت
\lr{A shengdebyte B}
می‌خوانیم! قبل از خواندن بقیه مسئله حتماً سه‌بار با خود بخوانید
\lr{A shengdebyte B}
و تکرار کنید چرا که کلید حل مسائل زیر در همین است!

\شروع{زیرمسائل}
\زیرمسئله[4]
الگوریتمی با 
$o(M^2)$
(یعنی بهتر از 
$M^2$
)
ارائه دهید که با داشتن $A$ و $B$ بتواند $A*B$ را خروجی دهد.

\زیرمسئله[8]
این مسئله کاربرد های زیادی به خصوص در مسئله کوله پشتی دارد. فرض کنید $n$ تا قفسه اشیا داریم که هر قفسه تعدادی شئ با وزن‌هایی بین $1$ تا $M$ دارد و فقسه‌ها عینا مانند هم هستند یعنی اشیا موجود در قفسه‌های اول و دوم و سوم و ... همگی برابر است. شما می‌خواهید از هر قفسه حداکثر یک شئ بردارید و می‌خواهید ببینید چه وزن‌هایی از بین $1$ تا $M$ را می‌توانید با خود حمل کنید. با استفاده از نتیجه بخش قبل الگوریتمی از
$o(M^2 log_2n)$
برای اینکار ارائه دهید.
\پایان{زیرمسائل}

\begin{itemize}
	\item[راهنمایی: ]
	یرای بخش الف یک عدد با $M+1$ رقم در مبنای $M+1$ متناظر هر مجموعه کنید که ارقامش یا صفر هستند و یا ۱ سپس از یکی از الگوریتم‌های ضرب اعداد بزرگ استفاده کنید تا نتیجه را به‌دست آورید.
\end{itemize}


%\متن‌سیاه{سوال عملی اول:}\\
%\متن‌سیاه{دراپ‌تست}\\

%\متن‌سیاه{ورودی:}\\
%در خط اول ورودی، ابتدا k و سپس m داده می‌شود.\\
%\متن‌سیاه{خروجی:}\\
%در تنها سطر خروجی، تعداد آیفون‌های مورد نیاز می‌آید. اگر در هیچ صورتی چنین کاری ممکن نبود ۱- چاپ می‌شود.\\
%\متن‌سیاه{محدودیت‌ها:}\\
%$1\leq k\leq 10^4$

%$m\leq10^9$
%\\
%\textbf{\متن‌سیاه{ورودی و خروجی نمونه}\\
%% تست اول
%\begin{latin}
% \begin{tabularx}{\textwidth}{|X | X|}
% \hline
% \textbf{stdin} &
% \textbf{stdout}\\
% \hline
% 100 100&
% 1\\
% \hline
% 
% \end{tabularx}
%\end{latin}
%
%% تست دوم
%
%\begin{latin}
% \begin{tabularx}{\textwidth}{|X | X|}
% \hline
% \textbf{stdin} &
% \textbf{stdout}\\
% \hline
% 340 20&
% 3\\
% \hline
% 
% \end{tabularx}
%\end{latin}}
% صورت اصلی مسئله را در این‌جا بنویسید.



\solution{
	این سؤال امتیازی شده‌است!
	\begin{itemize}
	
	
	\item[الف)]
		
	هر مجموعه را با عددی $M+1$ رقمی در مبنای $M+1$ جایگزین می‌کنیم که رقم $i$ام آن برابر ۱ است اگر و تنها اگر عضو $i$ در مجموعه باشد (رقم صفرم را هم برابر ۱ می‌گذاریم) حالا در صورت ضرب کردن دو عدد در هم رقمی مثل
	$j$
	بیشتر از صفر است اگر و تنها اگر در $A*B$ باشد چرا که:
	$$num(A) \times num(B) = \underset{0 \le d \le M+1}{\sum} (\underset{0 \le i \le d}{\sum} 10^iA_i\times 10^{d-i}B_i)$$
	$$=num(A) \times num(B) = \underset{0 \le d \le M+1}{\sum} 10^d (\underset{0 \le i \le d} {\sum} A_iB_i)$$
	چون رقم‌های $A$ و $B$ صفر یا یکند اگر مقدار جمع داخلی بیشتر از صفر باشد یعنی دو رقم داشتیم که جمع جایگاه‌هایشان برابر $d$شده‌است و چون $(M+1)$ بر یک نداریم این موضوع دقیقا در رقم $d$ ام 
	$num(A) \times num(B)$
	خودش را نشان می‌دهد. پس کافیست این دو را ضرب کنیم و ببینیم کدام رقم‌ها ناصفرند. برای ضرب اعداد بزرگ اگر از الگوریتم‌هایی مثل کاراتسوبا استفاده کنیم می‌توانیم به زمانی بهتر از 
	$M^2$
	برسیم. توجه کنید عدد حاصل ممکن است بیشتر از $M+1$ رقم داشته باشد ولی باید صرفا آن رقم‌های اضافی را نادیده بگیریم.
	
	\item [ب)]
	کاری که کافیست بکنیم این است که عدد متناظر را به توان $n$ برسانیم با استفاده از روش تقسیم و حل و استفاده از روش بخش الف مثل کاراتسوبا می‌توانیم اینکار را در
	$\mathcal{O}(M^{log_23} log n)$
	انجام دهیم. ولی باید حواسمان باشد پس از هر ضرب رقم‌های بیشتر از صفر را برابر ۱ کنیم تا
	$(M+1)$
	بر یک بوجود نیاید.
\end{itemize}

}

%%%%%%%%%%%%%%%%%%%%%%%%%%%%%%%%%%%%%%%%%%%%%%%%%%%%%
% برخی دستورات مفید برای تمرین‌ها. %
%%%%%%%%%%%%%%%%%%%%%%%%%%%%%%%%%%%%%%%%%%%%%%%%%%%%%

% فرمول‌های نمونه:
%$f(x) = \Theta(x)$
%$T(x, y) \leq \log(x) + 2^y + \binom{2n}{n}$
% {\tt code\_function}


% لیست موارد بدون شماره‌گذاری:
%\شروع{فقرات}
% \فقره مورد اول 
% \شروع{فقرات}
% \فقره مورد اول داخلی 
% \شروع{فقرات}
% \فقره اولین مورد خیلی داخلی 
% \فقره دومین مورد خیلی داخلی 
% \پایان{فقرات}
% \فقره مورد دوم داخلی 
% \پایان{فقرات}
% \فقره مورد دوم 
%\پایان{فقرات}

% لیست موارد شماره‌گذاری شده:
%\شروع{شمارش}
% \فقره یک مورد 
% \فقره یک مورد دیگر 
% \فقره باز هم یک مورد دیگر
%\پایان{شمارش}

% برای نوشتن یک فرمول چند خطی که یک شماره بخوره. برای حالت‌های دیگر به سایت زیر مراجعه کنید:
% https://www.overleaf.com/learn/latex/Aligning_equations_with_amsmath
%\begin{align} 
% \begin{split}
% (x+y)^3 &= (x+y)^2(x+y) \\
% &= (x^2+2xy+y^2)(x+y) \\
% &= (x^3+2x^2y+xy^2) + (x^2y+2xy^2+y^3) \\
% &= x^3+3x^2y+3xy^2+y^3
% \end{split} 
%\end{align}

% نوشتن ماتریس و آرایه دوبعدی
%\begin{align}
% A = 
% \begin{bmatrix}
% A_{11} & A_{21} \\
% A_{21} & A_{22}
% \end{bmatrix}
%\end{align}

% اضافه کردن تصویر با شرح. برای اضافه کردن تصویر بدون شرح، فقط دستور درج‌تصویر را استفاده کنید.
%\شروع{شکل}[h!]
%\تنظیم‌ازوسط
% \درج‌تصویر[پهنا=0.5\پهنای‌متن]{img.jpg}
%\شرح{شرح تصویر}
%\برچسب{تصویر:۱}
%\پایان{شکل}
%\پایان{زیرمسائل}

%کد زیر را ببینید.
%%\lstinputlisting{source_filename.py}
%\begin{latin}
%\begin{lstlisting}
%# Import the modules
%import sys
%import random
%
%ans = True
%
%while ans:
%question = raw_input("Ask the magic 8 ball a question: (press enter to quit) ")
%
%answers = random.randint(1,8)
%
%if question == "":
%sys.exit()
%
%elif answers == 1:
%print "It is certain"
%
%elif answers == 2:
%print "Outlook good"
%
%elif answers == 3:
%print "You may rely on it"
%
%elif answers == 4:
%print "Ask again later"
%
%elif answers == 5:
%print "Concentrate and ask again"
%
%elif answers == 6:
%print "Reply hazy, try again"
%
%elif answers == 7:
%print "My reply is no"
%
%elif answers == 8:
%print "My sources say no"
%\end{lstlisting}
%\end{latin}